\documentclass[a4paper,12pt]{report}

\usepackage[utf8]{inputenc} % vagy latin2 helyett utf8
\usepackage[T1]{fontenc}      % karakterkódolás
%\usepackage[magyar]{babel}    % magyar beállítások
\frenchspacing                % helyközök
%\usepackage{times}           % betûtípus
\usepackage{lmodern}          %   vagy inkább ez

\usepackage[margin=2.5cm,left=3.5cm,includeheadfoot]{geometry}
                              % margók
\usepackage{graphicx}         % képekhez
\usepackage{setspace}         % sorköz
\onehalfspacing               % másfeles




\begin{document}

% ------------------------------------------------------------------------------
% Címlap

\begin{titlepage}

\noindent
\parbox[m]{0.2\textwidth}{
    \includegraphics[width=0.2\textwidth]{img/elte_logo_colored.eps}
}
\hfill
\parbox[m]{0.7\textwidth}{
    \begin{center}
    \begin{large}
    \textsc{
        Eötvös Loránd University\\
        \vspace{0.5pc}
        Faculty of Informatics\\
        \vspace{0.5pc}
        Department of Programming Languages and Compilers\\
    }
    \end{large}
    \end{center}
}

\vspace{1pc}
\hrule

\vfill

\begin{center}
    {\LARGE Tool support for C++ template metaprogramming}
\end{center}

\vfill

\noindent
\hspace*{0.05\textwidth}
\parbox{0.45\textwidth}{
    {\it Supervisor:}
    \bigskip

    {\Large Zoltán Porkoláb}
    \smallskip

    Associate Professor
}
\hfill
\parbox{0.45\textwidth}{
    {\it Author:}
    \bigskip

    {\Large András János Kucsma}
    \smallskip

    Software Information Technology BSc
}


\vfill

\begin{center}
    {\large {\it Budapest, 2014}}
\end{center}

\end{titlepage}


% ------------------------------------------------------------------------------
% Témabejelentõ

\vspace*{\fill}
\begin{center}
Temabejelento
\end{center}
\vfill
\thispagestyle{empty}
\newpage
\setcounter{page}{1}

% ------------------------------------------------------------------------------
% Tartalomjegyzék

\tableofcontents

% ------------------------------------------------------------------------------


\chapter{Introduction}

A Bevezeto a temavalasztas indoklasat és a megoldando feladat rovid, kozertheto leirasat tartalmazza.


% ------------------------------------------------------------------------------


\chapter{Felhasználói dokumentáció}

A Felhasználói dokumentáció tartalmazza
\begin{itemize}
\item a megoldott probléma rövid megfogalmazását,
\item a felhasznált módszerek rövid leírását,
\item a program használatához szükséges összes információt.
\end{itemize}


% ------------------------------------------------------------------------------


\chapter{Fejlesztõi dokumentáció}

A Fejlesztõi dokumentáció tartlmazza
\begin{itemize}
\item a probléma részletes specifikációját,
\item a felhasznált módszerek részletes leírását, a használt fogalmak definícióját,
\item a program logikai és fizikai szerkezetének leírásár (adatszerkezetek, adatbázisok, modulfelbontás),
\item a tesztelési tervet és a tesztelés eredményeit.
\end{itemize}


% ------------------------------------------------------------------------------


% ------------------------------------------------------------------------------

\begin{thebibliography}{9}
\bibitem{v} Valaki, valami.
\end{thebibliography}

\end{document}
